\section{奖励塑形与时间信用分配}

\subsection{卷积 LSTM 和 reward shaping}

\begin{figure}[H]
\centering
\includegraphics[width=0.85\textwidth]{slide-24.jpg}
\caption{Lecture slide 24 描绘的 reward shaping 示例：在 multi-stage task 中用 intermediate reward guidance 缩短训练时间。}
\end{figure}

课程强调在 multi-stage/ hierarchical task 中提供 intermediate reward 能显著加速 critic 收敛。做法是把 reward 分解为 \textit{completion reward} + \textit{auxiliary reward}，并用 dynamic weighting 控制 gradient magnitudes。

\begin{knowledgebox}{Reward shaping 的 guardrails}
辅助 reward 必须保持 potential-based（即不改变 optimal policy），否则会鼓励 agent 偏离真正 objective。通常引入 shaping 函数 $F(s,a,s')$ 并确保 $\sum_t \gamma^t F(s_t,a_t,s_{t+1})$ 的总和为零。
\end{knowledgebox}

\subsection{可微分规划器与 credit assignment}

讲者提到一个实践 trick：在复杂控制中用 differentiable planner 估算未来 trajectory，再把 planner 的 soft value 作为 critic 监督信号。这样能显著减少 long-term credit assignment 的 drift。

\begin{warningbox}{Soft planner 需要额外计算}
引入 differentiable planning 会增加每步推理开销，必须确保该 planner 能在 low latency 环境里运行，否则会毁掉 real-time control。
\end{warningbox}

\subsection{本章小结}

Reward shaping + planner 信号是处理 multi-stage task 和 long-horizon credit assignment 的现实解决方案，但要确保 shaping 保持 potential-based，同时不抢占 inference budget。

